\documentclass[aspectratio=169,11pt]{beamer}
\usetheme{default}
\usefonttheme{professionalfonts}


\usepackage{graphicx}
\usepackage{booktabs}
\setbeamertemplate{navigation symbols}{}
\definecolor{ink}{HTML}{1F2937}
\definecolor{accent}{HTML}{9A3412}
\definecolor{muted}{HTML}{6B7280}
\setbeamercolor{frametitle}{fg=ink}
\setbeamercolor{title}{fg=ink}
\setbeamercolor{normal text}{fg=ink}
\setbeamercolor{itemize item}{fg=accent}
\setbeamercolor{itemize subitem}{fg=muted}
\setbeamerfont{frametitle}{size=\LARGE,series=\bfseries}
\setbeamerfont{title}{size=\Huge,series=\bfseries}
\setbeamerfont{subtitle}{size=\Large}
\setbeamercolor{subtitle}{fg=accent}
\setbeamertemplate{itemize item}{\textbullet}
\setbeamertemplate{itemize subitem}{--}
\setbeamertemplate{footline}{\hfill\color{muted}\scriptsize\insertframenumber\hspace{1.2em}\vspace{0.6em}}
\newcommand{\thread}[1]{{\color{accent}\bfseries #1}}
\newcommand{\credit}[1]{\par\vspace{0.3em}{\color{muted}\tiny #1}}
% \imgslide{title}{image}{image width fraction}{bullets}{credit}
\newcommand{\imgslide}[5]{\begin{frame}{#1}\begin{columns}[c]\begin{column}{#3\textwidth}\centering\includegraphics[width=\linewidth,height=0.78\textheight,keepaspectratio]{#2}\credit{#5}\end{column}\begin{column}{\dimexpr0.96\textwidth-#3\textwidth\relax}\large #4\end{column}\end{columns}\end{frame}}

\title{The Origins of Modern Computing}
\subtitle{CSCI 264 \textbullet{} Lecture 11 \textbullet{} Dive into Systems \S5.1}
\date{Monday, September 21, 2026}
\author{}

\begin{document}

\begin{frame}
\titlepage
\end{frame}

\begin{frame}{Listen for three threads}
\Large
\begin{itemize}\setlength\itemsep{1.2em}
\item \thread{Women in early computing} --- who did the work, and who got the credit
\item \thread{The military} --- who paid for the machines, and why
\item \thread{John von Neumann} --- and the document that put his name on an architecture
\end{itemize}
\end{frame}

\begin{frame}{A timeline}
\Large
\centering
\begin{tabular}{r@{\hspace{1.5em}}l}
\toprule
1837 & Babbage \& Lovelace: the Analytical Engine \\
1937 & Shannon's thesis; Turing's machine \\
1943 & Colossus at Bletchley Park \\
1944 & Harvard Mark I \\
1945 & ENIAC completed; the \emph{First Draft} circulated \\
1946 & ENIAC unveiled; Turing's ACE report \\
\bottomrule
\end{tabular}
\end{frame}

\begin{frame}{``Computer'' was a job title}
\Large
\begin{itemize}\setlength\itemsep{1em}
\item 1700s--early 1900s: a \emph{computer} was \emph{one who computes}
\item A person doing calculations by hand, usually one of a large team
\item Most often a woman
\item Every machine today was built to replace a room full of these people
\end{itemize}
\end{frame}

\imgslide{Babbage and Lovelace}{figures/ada-lovelace.jpg}{0.34}{\begin{itemize}\setlength\itemsep{0.9em}\item The \textbf{Analytical Engine}: a mechanical, general-purpose design\item Ada Lovelace published the first algorithm meant for a machine\item ``First computer'' is \emph{retrospective}\item Later designers mostly didn't know their work\end{itemize}}{Ada Lovelace, watercolour by A.\,E.\ Chalon. Wikimedia Commons.}

\imgslide{A primordial soup of ideas}{figures/alan-turing.jpg}{0.32}{\begin{itemize}\setlength\itemsep{0.9em}\item \textbf{1937, Claude Shannon}: Boolean logic can be built from switches\item[] \quad $\to$ the basis of all digital circuits\item \textbf{1937, Alan Turing}: the Turing machine\item[] \quad $\to$ what a computer can and cannot do\item Neither of them built a machine\end{itemize}}{Alan Turing. Wikimedia Commons.}

\imgslide{Colossus and the Wrens}{figures/colossus-wrens.jpg}{0.55}{\begin{itemize}\setlength\itemsep{0.8em}\item 1943, Tommy Flowers, Bletchley Park\item Arguably the first programmable, digital, fully electronic computer\item Operated by the \thread{Wrens}\item Noted for cryptographic ability \ldots\item \ldots\ never made cryptographers\end{itemize}}{Colossus at Bletchley Park, 1943. Wikimedia Commons.}

\begin{frame}{Meanwhile, in the United States}
\Large
\begin{itemize}\setlength\itemsep{1em}
\item \textbf{Harvard Mark I} --- Howard Aiken, a U.S.\ Navy Reserve commander
\item \textbf{ENIAC} --- Eckert and Mauchly, University of Pennsylvania, 1945
\begin{itemize}\large
\item Digital, fully electronic, programmable, general purpose
\item Decimal, not binary
\end{itemize}
\item Built for \thread{U.S.\ Army Ordnance}
\item Purpose: \textbf{artillery firing tables}
\end{itemize}
\end{frame}

\begin{frame}{The people ENIAC was built to replace}
\Large
\begin{itemize}\setlength\itemsep{1em}
\item $\approx$\,200 women computing trajectories by hand at Penn
\item Slide rules, desk calculators, a differential analyzer
\item One trajectory: 20 minutes to several days
\item Kay McNulty's appointment: \textbf{``subprofessional''} grade
\end{itemize}
\credit{J.\,S.\ Light, ``When Computers Were Women,'' \emph{Technology and Culture} 40:3 (1999).}
\end{frame}

\begin{frame}{The first programmers}
\Large
\begin{itemize}\setlength\itemsep{1.2em}
\item The human computers became the first programmers
\item Because programming was considered \textbf{clerical}
\item The machine was the achievement; instructing it was typing
\end{itemize}
\end{frame}

\imgslide{Grace Hopper}{figures/grace-hopper.jpg}{0.44}{\begin{itemize}\setlength\itemsep{0.8em}\item Programmer on the Mark I; wrote its manual\item First compiler: \textbf{A-0}\item \textbf{FLOW-MATIC} $\to$ COBOL\item A commissioned officer in the \thread{U.S.\ Naval Reserve}\end{itemize}}{Cropper, Krishnan, Hopper, Rothberg at a UNIVAC I console, 1957. Computer History Museum 102741216.}

\imgslide{The ``first bug''}{figures/first-computer-bug-1947.jpg}{0.52}{\begin{itemize}\setlength\itemsep{0.9em}\item Harvard Mark II log, 1947\item A moth, taped in\item Not the first use of ``bug''\item Probably not written by Hopper\end{itemize}}{Harvard Mark II logbook, 1947. Wikimedia Commons.}

\imgslide{The ENIAC Six}{figures/eniac-programmers.jpg}{0.55}{\normalsize Jean Jennings Bartik\\Betty Snyder Holberton\\Kay McNulty\\Frances Bilas Spence\\Marlyn Wescoff Meltzer\\Ruth Lichterman Teitelbaum\par\vspace{0.8em}\begin{itemize}\item Given wiring diagrams, not a manual\item Invented flow charts, subroutines, nesting\end{itemize}}{ENIAC at the Moore School. U.S.\ Army photo, Wikimedia Commons.}

\begin{frame}{The credit did not follow}
\Large
\begin{itemize}\setlength\itemsep{1em}
\item February 15, 1946: ENIAC unveiled, running Jennings and Snyder's program
\item Congratulations went to Eckert and Mauchly
\item \emph{New York Times}: the work of ``100 trained men''
\item Photo captions: ``setting switches,'' ``plugging cables''
\item Reused as an Army recruiting ad --- with the women cropped out
\end{itemize}
\end{frame}

\begin{frame}{ENIAC's first real job}
\Huge
\centering
\vfill
Not a firing table.\\[0.8em]
\Large A feasibility calculation for a \thread{thermonuclear weapon}.
\vfill
\end{frame}

\imgslide{John von Neumann}{figures/von-neumann.jpg}{0.30}{\begin{itemize}\setlength\itemsep{0.9em}\item Budapest, 1903; Institute for Advanced Study, 1933\item Mentor to Turing at Princeton\item Manhattan Project $\to$ Mark I, ENIAC\item ENIAC's first run was his calculation\end{itemize}}{Von Neumann's grave, Princeton Cemetery. Photo: M.\ West.}

\begin{frame}{The \emph{First Draft of a Report on the EDVAC}}
\Large
\begin{itemize}\setlength\itemsep{1em}
\item June 1945
\item EDVAC: ENIAC's successor, proposed by Eckert and Mauchly
\item \textbf{Binary}, not decimal
\item \textbf{Stored program}: instructions and data both live in memory
\item[] \quad Instructions are just more bits
\end{itemize}
\end{frame}

\begin{frame}{One name on the cover}
\Large
\begin{itemize}\setlength\itemsep{1em}
\item A draft, unfinished --- blanks where the references should go
\item Months of group work at the Moore School
\item Circulated by \thread{Herman Goldstine, the Army's liaison}
\item Eckert and Mauchly furious: circulation undermined their patent claims
\item Turing's ACE report, 1946 --- but von Neumann's came first
\end{itemize}
\end{frame}

\begin{frame}{Afterword}
\Large
\begin{itemize}\setlength\itemsep{1.4em}
\item \textbf{1952} --- Alan Turing convicted of ``gross indecency''\\ \large Clearance revoked; forced hormone treatment; dead in 1954
\item \textbf{1997} --- the six ENIAC programmers inducted into the WITI Hall of Fame\\ \large Fifty-one years later; Ruth Lichterman Teitelbaum had died in 1986
\end{itemize}
\end{frame}

\begin{frame}{Two minutes, on your own}
\Huge
\centering
\vfill
Of everything today,\\ what surprised you most --- and why?
\vfill
\large\color{muted} You won't hand it in.
\end{frame}

\begin{frame}{Discussion}
\LARGE
Programming was given to women because it was the ``clerical'' half of the work. Today it is the high-status half.\par\vspace{1.2em}
\thread{What changed? What does that tell us about how a field decides which work is valuable?}
\end{frame}

\begin{frame}{Discussion}
\LARGE
Almost every machine today was built for a war, with an army's money, to solve an army's problem.\par\vspace{1.2em}
\thread{Do we get this architecture without WWII --- or just later, and differently?}
\end{frame}

\begin{frame}{Discussion}
\LARGE
Eckert and Mauchly built it. Six women programmed it. Von Neumann's name is on it.\par\vspace{1.2em}
\thread{Is ``von Neumann architecture'' the wrong name? Does it matter?}
\end{frame}

\begin{frame}{Discussion}
\LARGE
The Wrens were recorded as having real cryptographic ability, and stayed operators. Nobody wrote down a decision to exclude them.\par\vspace{1.2em}
\thread{So what would have had to be different?}
\end{frame}

\begin{frame}{Further reading}
\large
\begin{itemize}\setlength\itemsep{0.8em}
\item Jennifer S.\ Light, ``When Computers Were Women,'' \emph{Technology and Culture} 40:3 (1999), 455--483
\item Janet Abbate, \emph{Recoding Gender}
\item Kathy Kleiman, \emph{Proving Ground} and the documentary \emph{The Computers}
\item LeAnn Erickson, \emph{Top Secret Rosies} (PBS)
\end{itemize}
\vfill
{\color{muted}\normalsize Wednesday: we open the box --- the five units of the von Neumann architecture.}
\end{frame}

\end{document}
